\section{Three connected research questions}
Tax policy changes reporting behavior, and those reports inform later policy. This loop combines the efficiency--equity problem in optimal taxation~\cite{mirrlees1971,saez2001,nobel1996} with a new governance risk: opaque learning may undermine legitimacy.

\textbf{Q1---Economics:} Under asymmetric information, when can transparent, bounded policy updates sustain compliance and welfare better than opaque updates?\par
\textbf{Q2---Computation:} Can an auditable finite-state simulator compare those regimes while enforcing bounds and reproducing every transition?\par
\textbf{Q3---Behavior:} Is any cooperation explained by trust in legitimate authority, or only by payoff-based strategic learning?\par
Incentives define the benchmark, computation makes the loop checkable, and behavior determines whether transparency changes choices. Figure~\ref{fig:teaser} joins the three channels.

\section{Answer to Q1: incentives and intellectual lineage}
My hypothesis is that bounded explanations preserve cooperation only when trust or continuation value supplements material payoffs. For eight simultaneous rounds, a taxpayer chooses Comply (C) or Avoid (A), and an authority Transparent (T) or Opaque (O). Model-set payoffs are C/T $=(3,3)$, C/O $=(0,4)$, A/T $=(4,0)$, and A/O $=(1,1)$. Both see the rule, trust score, and prior log, but not private motives. A and O are strictly dominant, making A/O the unique one-shot Nash equilibrium~\cite{nash1950}; C/T is Pareto-superior. Finite repetition with unchanged payoffs does not reverse that benchmark, so persistent C/T requires future incentives, beliefs, or social preferences---not a displayed trust score alone.

\section{Answer to Q2: method and evaluation}
The client-side finite-state game adds payoffs, updates trust by $+10,-15,-10,-10$, clips it to $[0,100]$, and logs every transition. Verified demonstrations produced $24/24$/100 (eight C/T), $7/11$/0 (one C/O then seven A/O), and $23/19$/90 (mixed recovery), where entries are taxpayer payoff/authority payoff/trust. These checks validate code, not behavior. The smallest experiment randomizes explained versus opaque updates with fixed payoffs; a modification varies observability. Colab must reproduce the rules, while Hugging Face provides the interactive interface.

\clearpage\balance
\section{Answer to Q3: behavior and competing explanations}
The slippery-slope mechanism predicts that transparent authority builds trust and voluntary compliance~\cite{kirchler2008}; strategic learning instead predicts cooperation only when future rewards justify it. A preregistered $2\times2$ test can vary explanation and continuation probability, measuring choices, comprehension, legitimacy, and payoffs. An explanation effect surviving payoff and comprehension controls favors trust; disappearance without continuation incentives favors learning. The model would add either trust-dependent utility or belief-dependent continuation value and log them separately. A classroom sample cannot identify population or policy effects. I would reject the trust claim if explanation changes comprehension but not behavior or legitimacy.


\section{Advanced development and future directions}
Mirrlees established how private information constrains tax incentives~\cite{mirrlees1971,nobel1996}; Sutton and Barto's Turing-recognized work established learning from feedback~\cite{acm2024}. The enduring question is how institutions secure fair cooperation without complete surveillance. Today's software can update rapidly and log decisions, but when recommendations remain incentive-compatible, understandable, and legitimate is unresolved. I contribute a bounded-update model, reproducible audit, and test separating trust from continuation incentives.

\begin{table}[h]\caption{From laureate foundations to a new interdisciplinary contribution.}
\begin{tabularx}{\columnwidth}{@{}p{1.6cm}Y@{}}\toprule
\textbf{Horizon}&\textbf{Milestone and evidence}\\\midrule
Foundation&Mirrlees: private-information tax incentives; Sutton--Barto: feedback learning.\\
PS1&Auditable game; three path checks.\\
Next study&Vary explanation and continuation value.\\
Longer term&Add diverse samples, legal rules, and privacy audits.\\
2056&Contestable adaptive institutions supported by field evidence.\\\bottomrule
\end{tabularx}\end{table}

\noindent\textbf{Expected 2056 Nobel Prize and Turing Award contribution statement.}
By 2056, I aspire to be recognized for designing adaptive tax institutions that people can understand and challenge. Building on Mirrlees and Sutton--Barto, I will integrate public economics, auditable computation, and behavioral evidence to determine when learning improves opportunity without sacrificing privacy or legitimacy. Evidence must progress from reproducible pilots to discriminating experiments, cross-jurisdiction validation, and accountable deployment.

\subsection*{Open Science Statement}
\textbf{GitHub:} \GitHubLink\\
\textbf{Google Colab:} \ColabLink\par
Synthetic inputs and client-side code reproduce the three paths. Before submission, add the final commit and verify a clean Colab ``Run all.'' Source-specific licenses remain in their repositories.

\subsection*{Statement of Contribution to the UN's SDGs}
\textbf{Hugging Face:} \DemoLink\\
The open browser game supports \textbf{SDG 4---Quality Education}~\cite{sdg4}: students choose both players' actions, inspect the log, and compare paths to learn repeated games, transparency, and trust. A pre/post prediction task could assess learning; improvement remains untested.

\subsection*{Submission milestones}
First draft: September 13, 2026, 11:00 P.M. Beijing time (UTC+8). Class review: September 14. Proposed: rebuttal September 16 and revised submission September 20, both at 11:00 P.M. Update Appendix~\ref{app:abstract} with each reflection and revision.
